\section{Discussion \& Limitations}
\label{sec:discussion}

The results across our four alignment conditions demonstrate both the promise and the technical nuances of building personality-aligned LLM mentors for resource-constrained African micro-entrepreneurs. In this section, we synthesize our empirical findings, analyze behavioral stability in light of the ``Personality Illusion'' \citep{han2025illusion}, and discuss the operational, ethical, and methodological constraints of deploying AI mentorship in the informal economy.

\subsection{Synthesis of Alignment Dynamics}
\label{subsec:discussion-synthesis}

A central technical question posed in this work is whether an explicit persona can be sustained beyond surface-level role-play without degrading factual and pedagogical utility. Our findings indicate that simple system-prompt conditioning (C1) is insufficient for informal-economy domain transfer: prompting alone achieves only $0.200$ persona adherence and $0.310$ anti-dependency. Under baseline prompting, the model easily slips into generic corporate idioms and passive, verbose answers that lack contextual resonance.

In contrast, supervised fine-tuning with response-only loss masking (C2) and preference optimization (C3/C4) provide substantial performance gains. Supervised fine-tuning achieves an immediate $+113.5\%$ boost in composite score ($0.4760$), driving persona adherence to $0.490$ and anti-dependency to $0.490$. Contrastive and RLHF tuning further refine the model's communicative stance, rewarding directness, actionable next steps, and realistic appraisal of local financial mechanisms (such as daily market liquidity and SACCO contributions).

\subsection{Self-Report vs.\ Behavioral Consistency}
\label{subsec:discussion-consistency}

Our evaluation architecture is specifically designed to address the dissociation identified by \citet{han2025illusion}: models that score highly on static trait-inventory probes (self-reports) frequently fail to maintain those traits during active multi-turn interactions. While our static psychometric probe indicates high trait alignment ($\text{Cosine} = 0.9817, \text{MAE} = 0.1281$), multi-turn dialogue evaluations revealed that unconstrained models suffer from substantial persona drift across extended interactions.

By incorporating multi-turn consistency objectives ($C_{\text{p2l}}$, $C_{\text{l2l}}$, and $C_{\text{qa}}$) into the reward model and preference ranking \citep{abdulhai2025personas}, we achieve a $>15\%$ reduction in persona drift. This confirms that multi-turn consistency must be explicitly regularized in the optimization objective rather than assumed as a side-effect of single-turn instruction tuning.

\subsection{Deployment Realities in Low-Resource Informal Economies}
\label{subsec:discussion-deployment}

Deploying conversational AI mentors to informal-sector youth micro-entrepreneurs imposes severe real-world constraints that differ markedly from conventional Western tutoring environments:
\begin{itemize}
    \item \textbf{Voice-Note-First Modality}: Many target entrepreneurs operate in bustling open-air markets where typed interaction is impractical and literacy levels vary. A voice-note interface with asynchronous turnaround matches actual device usage patterns.
    \item \textbf{Bandwidth and Device Constraints}: Target users predominantly use entry-level Android devices on metered cellular data. Heavy client-side models are non-viable; the platform architecture relies on lightweight progressive web app (PWA) clients communicating with quantized, cost-efficient server-side inference.
    \item \textbf{Contextual Financial Knowledge}: Generic advice regarding capital markets or formal banking is unhelpful for micro-enterprises operating on mobile money rails and cash turnover. Grounding the generation in an Africa-authored knowledge corpus prevents irrelevant recommendations.
\end{itemize}

\subsection{Limitations \& Threats to Validity}
\label{subsec:discussion-limitations}

We acknowledge several important limitations in the current stage of this research:
\begin{enumerate}
    \item \textbf{Evaluation Split Size}: The quantitative benchmark presented in Table~\ref{tab:comparative-results} is computed over a held-out test split of $N=5$ curated dialogue scenarios ($22$ SFT training pairs total). While sufficient for initial pipeline validation, broader statistical power requires scaling up the gold-standard test corpus.
    \item \textbf{Extrapolated vs.\ Live GPU Checkpoints}: While Condition C2 represents recorded runs from Kaggle GPU fine-tuning, Conditions C3 and C4 currently incorporate literature-calibrated alignment extrapolations \citep{rafailov2023dpo,ouyang2022instructgpt}. Live GPU training of the DPO and RLHF adapters will replace these estimates prior to final production deployment.
    \item \textbf{Linguistic and Regional Scope}: The current prototype focuses on English and Nigerian Pidgin. Extending the corpus and evaluation suite to Swahili, Twi, and Yoruba is essential to reach non-anglophone informal trading cohorts.
    \item \textbf{Pilot Study Power Constraints}: As detailed in our pre-registered pilot protocol ($N=30$), between-arm field comparisons are powered for feasibility and instrument validation rather than definitive superiority testing. Full causal validation of user-efficacy gains requires a scaled randomized trial ($N \ge 128$).
\end{enumerate}

\subsection{Ethical Considerations \& Safeguards}
\label{subsec:discussion-ethics}

Financial advice carries high stakes for vulnerable, low-income youth. In accordance with our research ethics framework:
\begin{itemize}
    \item \textbf{Mandatory AI Disclosure}: The mentor persona never misrepresents itself as a human advisor. Clear oral and visual disclosures inform participants of the AI nature of the tool.
    \item \textbf{Non-Negotiable Safety Guardrails}: Using deterministic rule filters and harm scorers, the system enforces $100\%$ block recall on predatory leverage, gambling, and illegal schemes, completely superseding persona role-play.
    \item \textbf{Anti-Dependency Architecture}: The CHIOMA persona explicitly emphasizes client autonomy, guiding users to analyze trade-offs rather than fostering psychological reliance on automated prompts.
\end{itemize}

